% GENERATED FILE. Edit canonical lessons, then run npm run generate:books.
% canonical-source-hash: fnv1a64:672e61e6b9fbd5de
% canonical-lessons: ML-C557-uttuka, ML-C557-sasikkuka, ML-C557-lalikkuka, ML-C557-talolikkuka, ML-C557-ottuceral

\chapter{Support and Scolding}
\label{ch:ml-support-and-scolding}
\hlchaptermodality{557}

\begin{chapteropening}
\textbf{By the end of this chapter:} \emph{I can say to hand-feed, to scold, to pamper, to cuddle and a get-together.}

Complete the last lesson of chapter 557: I can say to hand-feed, to scold, to pamper, to cuddle and a get-together.
\end{chapteropening}

\section[\texorpdfstring{\ml{ഊട്ടുക} (ūṭṭuka)}{ഊട്ടുക (ūṭṭuka)}]{\ml{ഊട്ടുക} (ūṭṭuka) — to hand-feed; to feed}

\label{lesson:ML-C557-uttuka}



\begin{quote}
Before the new one: say the Malayalam for care, then the Malayalam for consideration.
\end{quote}

\subsection*{You'll want to know: \ml{ഊട്ടുക}}
\textbf{\ml{ഊട്ടുക}} — \emph{ūṭṭuka} — \textquotedblleft{}to hand-feed; to feed\textquotedblright{}.

\begin{quote}
\ml{അമ്മ} \ml{കുഞ്ഞിനെ} \ml{ചോറ്} \ml{ഊട്ടി}.
\end{quote}

\begin{quote}
\emph{amma kuññine cōṟŭ ūṭṭi.}
\end{quote}

\begin{quote}
Mother fed the baby rice by hand.
\end{quote}

Past: \ml{ഊട്ടി}, \emph{ūṭṭi}.

Feeding a child or a guest by hand; the past is \ml{ഊട്ടി}.

\begin{cousinweb}
\textbf{\ml{ഊട്ടുക}} (\emph{ūṭṭuka}) is native Dravidian: DEDR 600 lists it beside \ml{ഉണ്ണുക}, \textquotedblleft{}to eat\textquotedblright{}, and \ml{ഊണ്}, \textquotedblleft{}a meal\textquotedblright{}, as the verb of making someone eat (DEDR 600).
\end{cousinweb}

\begin{grammarlens}[title={What you've built}]
One more word for describing what you lived and felt.
\end{grammarlens}

\subsection*{Your turn}
\begin{itemize}
\raggedright
  \item \emph{Say it:} \emph{ūṭṭuka}
  \item \emph{Say it:} \emph{ūṭṭuka}, once more
  \item \emph{Say it:} \emph{amma kuññine cōṟŭ ūṭṭi.} — Mother fed the baby rice by hand.
  \item \emph{Recall:} say the Malayalam for care, then the Malayalam for consideration, then say \emph{ūṭṭuka} again
\end{itemize}

\begin{tcolorbox}[breakable,colback=teal!4,colframe=teal!35!black,title={Before you move on}]
What does \ml{ഊട്ടുക} mean? (\textquotedblleft{}To hand-feed; to feed\textquotedblright{}.) Say it once more.
\end{tcolorbox}
\section[\texorpdfstring{\ml{ശാസിക്കുക} (śāsikkuka)}{ശാസിക്കുക (śāsikkuka)}]{\ml{ശാസിക്കുക} (śāsikkuka) — to scold; to reprimand}

\label{lesson:ML-C557-sasikkuka}



\begin{quote}
Before the new one: say the Malayalam for consideration, then the Malayalam for to hand-feed.
\end{quote}

\subsection*{You'll want to know: \ml{ശാസിക്കുക}}
\textbf{\ml{ശാസിക്കുക}} — \emph{śāsikkuka} — \textquotedblleft{}to scold; to reprimand\textquotedblright{}.

\begin{quote}
\ml{ക്ലാസിൽ} \ml{സംസാരിച്ചതിന്} \ml{അധ്യാപകൻ} \ml{കുട്ടികളെ} \ml{ശാസിച്ചു}.
\end{quote}

\begin{quote}
\emph{klāsil saṁsāriccatinŭ adhyāpakan kuṭṭikaḷe śāsiccu.}
\end{quote}

\begin{quote}
The teacher scolded the children for talking in class.
\end{quote}

Past: \ml{ശാസിച്ചു}, \emph{śāsiccu}.

Firmer than \ml{വഴക്ക്} \ml{പറയുക}; the noun is \ml{ശാസന}.

\begin{cousinweb}
\textbf{\ml{ശാസിക്കുക}} (\emph{śāsikkuka}) is built from Sanskrit \emph{śāsana}, \textquotedblleft{}instructing, a command, punishment\textquotedblright{} (CDIAL 12420), with the Malayalam verb ending -\ml{ിക്കുക}.
\end{cousinweb}

\begin{grammarlens}[title={What you've built}]
One more word for describing what you lived and felt.
\end{grammarlens}

\subsection*{Your turn}
\begin{itemize}
\raggedright
  \item \emph{Say it:} \emph{śāsikkuka}
  \item \emph{Say it:} \emph{śāsikkuka}, once more
  \item \emph{Say it:} \emph{klāsil saṁsāriccatinŭ adhyāpakan kuṭṭikaḷe śāsiccu.} — The teacher scolded the children for talking in class.
  \item \emph{Recall:} say the Malayalam for consideration, then the Malayalam for to hand-feed, then say \emph{śāsikkuka} again
\end{itemize}

\begin{tcolorbox}[breakable,colback=teal!4,colframe=teal!35!black,title={Before you move on}]
What does \ml{ശാസിക്കുക} mean? (\textquotedblleft{}To scold; to reprimand\textquotedblright{}.) Say it once more.
\end{tcolorbox}
\section[\texorpdfstring{\ml{ലാളിക്കുക} (lāḷikkuka)}{ലാളിക്കുക (lāḷikkuka)}]{\ml{ലാളിക്കുക} (lāḷikkuka) — to pamper; to spoil}

\label{lesson:ML-C557-lalikkuka}



\begin{quote}
Before the new one: say the Malayalam for to hand-feed, then the Malayalam for to scold.
\end{quote}

\subsection*{You'll want to know: \ml{ലാളിക്കുക}}
\textbf{\ml{ലാളിക്കുക}} — \emph{lāḷikkuka} — \textquotedblleft{}to pamper; to spoil\textquotedblright{}.

\begin{quote}
\ml{മുത്തശ്ശി} \ml{എപ്പോഴും} \ml{ഇളയ} \ml{പേരക്കുട്ടിയെ} \ml{ലാളിക്കും}.
\end{quote}

\begin{quote}
\emph{muttaśśi eppōḻuṁ iḷaya pērakkuṭṭiye lāḷikkuṁ.}
\end{quote}

\begin{quote}
Grandmother always pampers her youngest grandchild.
\end{quote}

Past: \ml{ലാളിച്ചു}, \emph{lāḷiccu}.

Also to spoil with too much love.

\begin{cousinweb}
\textbf{\ml{ലാളിക്കുക}} (\emph{lāḷikkuka}) is built from Sanskrit \emph{lālana}, \textquotedblleft{}caressing\textquotedblright{} (CDIAL 11025), with the Malayalam verb ending -\ml{ിക്കുക}.
\end{cousinweb}

\begin{grammarlens}[title={What you've built}]
One more word for describing what you lived and felt.
\end{grammarlens}

\subsection*{Your turn}
\begin{itemize}
\raggedright
  \item \emph{Say it:} \emph{lāḷikkuka}
  \item \emph{Say it:} \emph{lāḷikkuka}, once more
  \item \emph{Say it:} \emph{muttaśśi eppōḻuṁ iḷaya pērakkuṭṭiye lāḷikkuṁ.} — Grandmother always pampers her youngest grandchild.
  \item \emph{Recall:} say the Malayalam for to hand-feed, then the Malayalam for to scold, then say \emph{lāḷikkuka} again
\end{itemize}

\begin{tcolorbox}[breakable,colback=teal!4,colframe=teal!35!black,title={Before you move on}]
What does \ml{ലാളിക്കുക} mean? (\textquotedblleft{}To pamper; to spoil\textquotedblright{}.) Say it once more.
\end{tcolorbox}
\section[\texorpdfstring{\ml{താലോലിക്കുക} (tālōlikkuka)}{താലോലിക്കുക (tālōlikkuka)}]{\ml{താലോലിക്കുക} (tālōlikkuka) — to cuddle; to fondle}

\label{lesson:ML-C557-talolikkuka}



\begin{quote}
Before the new one: say the Malayalam for to scold, then the Malayalam for to pamper.
\end{quote}

\subsection*{You'll want to know: \ml{താലോലിക്കുക}}
\textbf{\ml{താലോലിക്കുക}} — \emph{tālōlikkuka} — \textquotedblleft{}to cuddle; to fondle\textquotedblright{}.

\begin{quote}
\ml{അമ്മ} \ml{കുഞ്ഞിനെ} \ml{മടിയിൽ} \ml{കിടത്തി} \ml{താലോലിച്ചു}.
\end{quote}

\begin{quote}
\emph{amma kuññine maṭiyil kiṭatti tālōliccu.}
\end{quote}

\begin{quote}
Mother laid the baby in her lap and cuddled it.
\end{quote}

Past: \ml{താലോലിച്ചു}, \emph{tālōliccu}.

Fond touching and rocking of a small child.

\begin{cousinweb}
\textbf{\ml{താലോലിക്കുക}} (\emph{tālōlikkuka}) is a fond, spoken verb; no source we can open explains it, so its origin is not settled.
\end{cousinweb}

\begin{grammarlens}[title={What you've built}]
One more word for describing what you lived and felt.
\end{grammarlens}

\subsection*{Your turn}
\begin{itemize}
\raggedright
  \item \emph{Say it:} \emph{tālōlikkuka}
  \item \emph{Say it:} \emph{tālōlikkuka}, once more
  \item \emph{Say it:} \emph{amma kuññine maṭiyil kiṭatti tālōliccu.} — Mother laid the baby in her lap and cuddled it.
  \item \emph{Recall:} say the Malayalam for to scold, then the Malayalam for to pamper, then say \emph{tālōlikkuka} again
\end{itemize}

\begin{tcolorbox}[breakable,colback=teal!4,colframe=teal!35!black,title={Before you move on}]
What does \ml{താലോലിക്കുക} mean? (\textquotedblleft{}To cuddle; to fondle\textquotedblright{}.) Say it once more.
\end{tcolorbox}
\section[\texorpdfstring{\ml{ഒത്തുചേരൽ} (ottucēral)}{ഒത്തുചേരൽ (ottucēral)}]{\ml{ഒത്തുചേരൽ} (ottucēral) — a get-together; a reunion}

\label{lesson:ML-C557-ottuceral}



\begin{quote}
Before the new one: say the Malayalam for to pamper, then the Malayalam for to cuddle.
\end{quote}

\subsection*{You'll want to know: \ml{ഒത്തുചേരൽ}}
\textbf{\ml{ഒത്തുചേരൽ}} — \emph{ottucēral} — \textquotedblleft{}a get-together; a reunion\textquotedblright{}.

\begin{quote}
\ml{അവധിക്ക്} \ml{കുടുംബത്തിന്റെ} \ml{ഒരു} \ml{ഒത്തുചേരൽ} \ml{ഉണ്ടായിരുന്നു}.
\end{quote}

\begin{quote}
\emph{avadhikkŭ kuṭumbattinṟe oru ottucēral uṇṭāyirunnu.}
\end{quote}

\begin{quote}
There was a family get-together in the holidays.
\end{quote}

\ml{ഒത്ത്} together + \ml{ചേരൽ} a coming together, from \ml{ചേരുക}, to join.

\begin{cousinweb}
\textbf{\ml{ഒത്തുചേരൽ}} (\emph{ottucēral}) is native Dravidian, a compound of \ml{ഒത്ത്}, from \ml{ഒക്കുക}, \textquotedblleft{}to agree, to be alike\textquotedblright{} (DEDR 924), and \ml{ചേരൽ}, \textquotedblleft{}a joining\textquotedblright{}, from \ml{ചേരുക} (DEDR 2814).
\end{cousinweb}

\begin{grammarlens}[title={What you've built}]
One more word for describing what you lived and felt.
\end{grammarlens}

\subsection*{Your turn}
\begin{itemize}
\raggedright
  \item \emph{Say it:} \emph{ottucēral}
  \item \emph{Say it:} \emph{ottucēral}, once more
  \item \emph{Say it:} \emph{avadhikkŭ kuṭumbattinṟe oru ottucēral uṇṭāyirunnu.} — There was a family get-together in the holidays.
  \item \emph{Recall:} say the Malayalam for to pamper, then the Malayalam for to cuddle, then say \emph{ottucēral} again
\end{itemize}

\begin{tcolorbox}[breakable,colback=teal!4,colframe=teal!35!black,title={Before you move on}]
What does \ml{ഒത്തുചേരൽ} mean? (\textquotedblleft{}A get-together; a reunion\textquotedblright{}.) Say it once more.
\end{tcolorbox}
